%! AP Statistics — Unit 2, Lesson 2.3 — Guided Notes (Answer Key)
\documentclass[10pt]{article}
\usepackage{apstats-article}
\usepackage{apstats-key}

%=============================================================================
\begin{document}

\pageheader{Unit 2, Lesson 2.3}{Guided Notes --- Answer Key}
\begin{tabularx}{\linewidth}{X X X}
Name:\;\ans{Key} & Date:\; & Period:\;
\end{tabularx}

\vspace{0.12in}

\begin{objectivebox}
\small By the end of this lesson, I will be able to\dots\\[4pt]
\ansline{compute the conditional proportion for each group in a two-way table,}\\[1pt]
\ansline{calculate and interpret the difference in proportions and relative risk,}\\[1pt]
\ansline{compare both statistics to assess the strength of an association.}
\end{objectivebox}

\vspace{0.1in}

\begin{vocabbox}
\termblanklong{Conditional proportion ($\hat{p}$)}
\hfill\ans{the proportion of successes within a given group: $\hat{p} = \frac{\text{successes}}{\text{group total}}$}

\termblanklong{Difference in proportions}
\hfill\ans{$\hat{p}_1 - \hat{p}_2$; additive comparison between two groups' rates}

\termblanklong{Relative risk (risk ratio)}
\hfill\ans{$\hat{p}_1 / \hat{p}_2$; multiplicative comparison; how many times as likely}

\termblanklong{Baseline group}
\hfill\ans{the reference group in the denominator of the relative risk}
\end{vocabbox}

\vspace{0.1in}

\begin{hookbox}
\small\textbf{Scenario:} 400 vaccinated, 32 flu cases; 400 unvaccinated, 80 flu cases.

\begin{enumerate}[leftmargin=1.5em, itemsep=2pt]
  \item Vaccinated: $\hat{p}_V = \ans{32/400 = 0.08}$ \quad Unvaccinated: $\hat{p}_U = \ans{80/400 = 0.20}$
  \item \ansline{We need a statistic --- either the difference (additive) or the ratio (multiplicative).}
\end{enumerate}
\end{hookbox}

\vspace{0.1in}

\begin{notesbox}{Step 1 --- Difference in Proportions (UNC-1.Q, UNC-1.R)}
\small

\[
  \hat{p}_V - \hat{p}_U = \frac{\ans{32}}{\ans{400}} - \frac{\ans{80}}{\ans{400}}
  = \ans{0.08} - \ans{0.20} = \ans{-0.12}
\]

\textbf{Interpretation:}\\[4pt]
\ansline{The proportion of vaccinated participants who got the flu (8\%) was 0.12 (12 percentage points) \emph{lower} than the proportion of unvaccinated participants who got the flu (20\%).}

\vspace{4pt}
\textbf{Key properties:}
\begin{itemize}[leftmargin=1.5em, itemsep=2pt]
  \item Value of \ans{0} $\to$ no difference.
  \item Positive $\to$ Group 1 rate is \ans{higher}; Negative $\to$ Group 1 rate is \ans{lower}.
  \item Units: \ans{proportion (or percentage points)}.
\end{itemize}
\end{notesbox}

\vspace{0.1in}

\begin{notesbox}{Step 2 --- Relative Risk (UNC-1.Q, UNC-1.R)}
\small

The denominator is the \textbf{\ans{baseline (unvaccinated)}} group.

\[
  \text{RR} = \frac{\hat{p}_V}{\hat{p}_U} = \frac{\ans{0.08}}{\ans{0.20}} = \ans{0.40}
\]

\textbf{Interpretation:}\\[4pt]
\ansline{Vaccinated participants were 0.40 times as likely to get the flu as unvaccinated
participants --- i.e., vaccination reduced the likelihood of flu to 40\% of the unvaccinated rate.}

\vspace{4pt}
\textbf{Key properties:}
\begin{itemize}[leftmargin=1.5em, itemsep=2pt]
  \item Value of \ans{1} $\to$ no difference.
  \item RR $> 1$ $\to$ Group 1 rate is \ans{higher}; RR $< 1$ $\to$ Group 1 rate is \ans{lower}.
  \item Not a percentage-point difference.
\end{itemize}
\end{notesbox}

\vspace{0.1in}

\begin{notesbox}{Step 3 --- When Do They Tell Different Stories?}
\small
\begin{multicols}{2}

\textbf{Context A: Rare outcome}\\
Difference in proportions: \ans{0.001}\\
Relative risk: \ans{2}

\vspace{4pt}
\ansline{RR is more striking --- ``twice as likely'' sounds alarming even though the absolute difference is tiny.}

\columnbreak

\textbf{Context B: Common outcome}\\
Difference in proportions: \ans{0.20}\\
Relative risk: \ans{1.5}

\vspace{4pt}
\ansline{The absolute difference (20 percentage points) is more striking here.}

\end{multicols}

\vspace{4pt}
\textbf{Takeaway:} A small \ans{absolute difference} can correspond to a large \ans{relative risk}.
Always report \ans{both statistics in context}.
\end{notesbox}

\vspace{0.1in}

\begin{spiralbox}
\small
\begin{multicols}{2}
\textbf{How this connects:} [same as student version]

\columnbreak

\ansline{Relative risk is preferred when the outcome is rare, because the multiplicative
comparison is more interpretable than a tiny absolute difference. In medical research, both
are routinely reported together.}
\end{multicols}
\end{spiralbox}

\begin{teachernote}
\textbf{Hook item 1:} Students should recall that $\hat{p} = \text{successes}/\text{group total}$
from Lesson 2.2. Circulate to check that they divide by 400, not by the grand total 800.

\textbf{Step 2 RR $< 1$:} Many students expect RR to always be $>1$. Emphasize that which group
is Group 1 is arbitrary --- swapping gives RR $= 1/0.40 = 2.5$ (unvaccinated are 2.5 times as
likely to get flu). Both are correct; direction must be stated.

\textbf{Step 3:} Context A is the classic ``scary headline'' situation (drug doubles cancer risk
from 0.001 to 0.002). Use to discuss media literacy and statistical communication.
\end{teachernote}

\end{document}
